\documentclass[11pt,a4paper]{article}

\usepackage[margin=2.5cm]{geometry}
\usepackage{amsmath,amssymb}
\usepackage[hidelinks]{hyperref}
\usepackage{microtype}

\input{glyphtounicode}
\pdfgentounicode=1

\newcommand{\SU}{\mathrm{SU}}
\newcommand{\SO}{\mathrm{SO}}
\newcommand{\C}{\mathbb{C}}
\newcommand{\R}{\mathbb{R}}
\newcommand{\Z}{\mathbb{Z}}
\newcommand{\Hcolor}{[\mathrm{H\text{-}color}]}

\title{\textbf{Withdrawal Notice for O31:\\
Unsupported Colour-Triplet Co-admissibility and Gauge-Group Claims}\\[0.5em]
\large Cosmochrony Spectral Admissibility Sub-programme}

\author{J\'er\^ome Beau\\[0.3em]
\small \texttt{jerome.beau@cosmochrony.org}}

\date{2026\\[2pt]\small Version 2.0}

\hypersetup{
  pdftitle={Withdrawal Notice for O31: Unsupported Colour-Triplet Co-admissibility and Gauge-Group Claims},
  pdfauthor={Jérôme Beau},
  pdfsubject={Withdrawal of unsupported O31 colour-triplet co-admissibility and gauge-group claims},
  pdfkeywords={Cosmochrony, withdrawal notice, colour triplet, co-admissibility, Heisenberg group,
    Weil representation, gauge group, counterexample}
}

\begin{document}

  \maketitle

  \begin{abstract}
    This version replaces the O31 paper in the same documentary chain.
    It is a withdrawal notice, not a corrected paper.
    Eight results of the previous version are withdrawn: the uniqueness of $\SU(3)$ among connected
    compact subgroups acting irreducibly on $\C^3$; the irreducibility of the group of
    Born--Infeld-compatible fibre-preserving transformations; the co-admissibility of the individual
    profiles $\sigma_{c_i}$; the arithmetic criterion that colour triplets exist if and only if
    $q \equiv 1 \pmod 3$; the metaplectic identity as used, which changes the central character
    while the automorphism invoked fixes the centre; the $\Delta(27)$ spectral proposition; the
    mistyped rank-eight test in $\mathrm{End}(\C^3)$; and every derivation of a gauge factor or of
    the Standard Model gauge group $\SU(3) \times \SU(2) \times \mathrm{U}(1)$.
    Three counterexamples are given in full.
    $\SO(3) \subset \SU(3)$ is a proper connected compact subgroup acting irreducibly on $\C^3$, so
    irreducibility with determinant one never singles out $\SU(3)$.
    The diagonal torus $T^2 \subset \SU(3)$ satisfies every constraint that was imposed while acting
    reducibly, and requiring each sector to be mapped to itself in fact forces reducibility.
    For every prime $q \geq 7$ the set $\{1, 2, q-3\}$ is an admissible colour triplet under the
    published definition, so the arithmetic criterion is false.
    The remaining scalar, arithmetic and discrete-group material is under re-audit and is not
    asserted.
    No new positive conclusion about groups, gauge factors, or the Standard Model gauge group is
    asserted here, and none should be inferred.
    The previous version remains available and citable in this record's version history, for
    traceability.
  \end{abstract}

  \paragraph{Keywords.}
    Cosmochrony; withdrawal notice; colour triplet; co-admissibility; Heisenberg group;
    Weil representation; gauge group; counterexample.

  \section*{Status}

    The results listed in Section~\ref{sec:withdrawn} are withdrawn; the counterexamples that
    require their withdrawal are given in Section~\ref{sec:counter}; and the remainder of the
    previous version is under re-audit (Section~\ref{sec:reaudit}).

  \section{Withdrawn results}
    \label{sec:withdrawn}

    The following are withdrawn.

    \begin{enumerate}
      \item \textbf{Uniqueness of $\SU(3)$.} The claim that the only connected compact subgroup of
      $\SU(3)$ acting unitarily and irreducibly on $\C^3$ is $\SU(3)$ itself.
      \item \textbf{Irreducibility.} The claim that the group of Born--Infeld-compatible
      fibre-preserving transformations of the colour eigenspace acts irreducibly.
      \item \textbf{Co-admissibility of the individual profiles.} The claim that the invariance of a
      simultaneously rescaled triplet block yields the equality of the three individual profiles
      $\sigma_{c_1}$, $\sigma_{c_2}$, $\sigma_{c_3}$.
      \item \textbf{The arithmetic criterion.} The claim that colour triplets exist if and only if
      $q \equiv 1 \pmod 3$.
      \item \textbf{The metaplectic identity as used.} The step that changes the central character
      while the automorphism invoked fixes the centre.
      \item \textbf{The $\Delta(27)$ spectral proposition} relating its characters and adjacency
      eigenvalues.
      \item \textbf{The rank-$8$ test} in $\mathrm{End}(\C^3)$, which is mistyped:
      $\dim_{\C} \mathrm{End}(\C^3) = 9$ and the stated test cannot be performed as written.
      \item \textbf{Every derivation of a gauge factor}, and of the Standard Model gauge group
      $\SU(3) \times \SU(2) \times \mathrm{U}(1)$ as a composite. The determinant-one condition, the
      commutation of the putative factors and the direct-product structure are likewise not
      established.
    \end{enumerate}

  \section{Counterexamples}
    \label{sec:counter}

    \paragraph{Uniqueness fails.}
      $\SO(3) \subset \SU(3)$ --- the real orthogonal matrices of determinant one, read as complex
      matrices --- is a proper connected compact subgroup, since
      $\dim \SO(3) = 3 < 8 = \dim \SU(3)$.
      Its action on $\C^3$ is the complexification of its defining representation, the
      $\ell = 1$ representation, whose commutant is one-dimensional; by Schur's lemma it is
      irreducible over $\C$.
      Irreducibility together with determinant one therefore never singles out $\SU(3)$.

    \paragraph{Irreducibility fails, and is in fact excluded by the stated constraints.}
      The diagonal torus
      \[
        T^2 = \mathopen{\{}\operatorname{diag}(e^{i\alpha}, e^{i\beta}, e^{-i(\alpha+\beta)})
        \;:\; \alpha, \beta \in \R \mathclose{\}} \subset \SU(3) .
      \]
      maps each coordinate line to itself, has determinant one, and maps no sector to a conjugate
      sector.  It satisfies every constraint that was imposed, while preserving the three coordinate
      lines --- that is, acting reducibly.
      More sharply: a requirement that each sector be mapped to itself makes any admissible
      transformation block-diagonal, hence reducible by construction.

    \paragraph{The arithmetic criterion fails.}
      For every prime $q \geq 7$, the set
      \[
        \{1,\, 2,\, q-3\}
      \]
      is an admissible colour triplet under the published definition.
      Its elements lie in $(\Z/q\Z)^{*}$ and are distinct precisely for $q \geq 7$; their sum is
      $q \equiv 0$; and their three pairwise sums are $3$, $q-2$ and $q-1$, none divisible by $q$,
      so pairwise non-conjugacy holds.
      The case $q = 5$ fails only because $q - 3 = 2$ duplicates an element.
      Since this construction is available for every prime $q \geq 7$ regardless of
      $q \bmod 3$, the criterion ``triplets exist if and only if $q \equiv 1 \pmod 3$'' is false.
      For instance $q = 11$ admits $\{1, 2, 8\}$, and $11 \equiv 2 \pmod 3$.

  \section{Under re-audit}
    \label{sec:reaudit}

    The scalar, arithmetic and discrete-group material of the previous version is under re-audit and
    is not asserted here.
    This includes the generator-twist isospectrality statement at fixed central character, which may
    be recoverable; if an autonomous proof is established, it will appear as a separate work rather
    than as a claim of this record.

    Nothing in the withdrawn list is replaced by a weaker positive statement in this notice.

  \section*{Scope}

    This notice concerns the O31 record only.
    It makes no statement about the status of any other paper of the programme, and any dependent
    claim elsewhere that relies on a withdrawn result above should be treated as unsupported until
    separately re-examined.

\end{document}
